\documentclass[11pt]{article}
\usepackage[margin=1in]{geometry}
\usepackage{titlesec}
\usepackage{enumitem}
\usepackage{nopageno}

% NSF-standard run-in bold headers
\titleformat{\section}[runin]{\normalfont\bfseries}{\thesection}{1em}{}[]
\titlespacing*{\section}{0pt}{0.6\baselineskip}{1ex}

\begin{document}

\noindent \textbf{Mentoring Plan}

\section*{Overview}
The Phase I CAMEL Collaborative Network (CAMEL-CN) involves graduate students from Penn State University who will be advised by the PI and Co-PIs across the fields of architectural engineering, materials science, and the learning sciences. The project aims to develop well-rounded researchers prepared to address complex challenges at the intersection of mathematics education research and data science. The success of this plan will be assessed through annual progress reports, student feedback, and professional achievements.

\section*{1. Regular One-on-One Meetings}
PI Napolitano and Co-PIs Yao will hold weekly meetings with their respective graduate students to discuss research progress, telemetry framework development, and pedagogical alignment. As applicable, PI/ other Co-PIs will serve on each other’s graduate committees to provide multifaceted support across data science and educational research domains.

\section*{2. Career Counseling and Professional Development}
Advisors will provide counseling to help students identify long-term goals in academia, industry, or educational policy. Students will be encouraged to present at major conferences such as ICLS. All students will have access to Penn State Graduate Career Services for mock interviewing and job search strategies.

\section*{3. Training in Grant Proposals and Scholarly Communication}
Students will receive training on preparing NSF grant proposals and peer-reviewed publications through the Penn State Graduate Writing Center. Advisors will review drafts of manuscripts regarding mathematical modeling development and the Data Translation Framework to refine writing skills. Students will also have access to PI Napolitano’s graduate-level research methods course materials.

\section*{4. Teaching and Mentoring Skill Development}
Graduate students will assist in developing modeling modules and delivering content during Summer Institutes. Advisors will provide guidance on effective techniques for facilitating ``messy" data experiences for high school learners. Students will be encouraged to pursue the Penn State Center for the Integration of Research, Teaching and Learning (CITRL) Certificate. Furthermore, graduate students will receive specific mentorship on collaborating with rural math teachers from CSATS, treating practitioners as intellectual co-researchers.

\section*{5. Interdisciplinary Collaboration and Equity Training}
Students will be trained to collaborate across disciplinary boundaries, including learning sciences and engineering. They will participate in university-led communication trainings such as ``Addressing Unconscious Bias as a Leader" and ``Workplace Diversity, Equity, Inclusion, and Belonging in Action" to prepare them for work with all types of American populations.

\section*{6. Responsible Professional Practices and Data Ethics}
All graduate students will receive training in the responsible conduct of research (RCR) through the PSU Scholarship and Research Integrity (SARI) program. Given the sensitive nature of process-level telemetry, students will receive specific training on secure data handling, PII scrubbing via pattern-matching algorithms, and the ethical framework for non-evaluative student monitoring. This ensures all project activities comply with Penn State IRB and FAIR data principles.

\end{document}